\section{Introduction}
\label{sec:introduction}


Neural scaling laws have ushered in new paradigms for interpreting deep learning dynamics. By establishing that performance gains follow predictable behavior relative to model size, dataset scale, and compute, stochastic optimization methods have been pushed toward increasingly extended training horizons while operating under tight efficiency constraints~\cite{hestness2017deep, kaplan2020scaling, henighan2020scaling, hoffmann2022chinchilla}. In this regime, optimization stability becomes an important practical requirement, since unstable training may result on slow convergence and critical loss spikes [add references here]. In practice, stability is often maintained through a set of heuristics, such as learning-rate decay, gradient clipping, and more recently Adam's second moments dynamic tuning~\cite{goyal2017accurate, smith2017cyclical, loshchilov2016sgdr, ioffe2015batch, krogh1991simple}[adjust the citations]. Yet, learning instabilities and oscilation can coexist with training loss decreases over longer intervals~\cite{cohen2021edge,damian2023self}. Reducing the learning rate or limiting gradient size can control fluctuations~\cite{smith2017cyclical,pascanu2013difficulty}, but a smoother loss curve need not mean faster loss reduction. The central question is, therefore, which changes should an learning intervention target to improve training?

Training becomes unstable when updates are too large relative to sharp directions of the loss landscape, so that they overshoot rather than converge towards a minimum~\cite{polyak1964methods, bertsekas1997nonlinear}. Stochastic extensions, such as the break-even point framework, argue that instability depends jointly on sharpness and curvature variance, particularly during an early phase in which learning rate and batch size strongly shape further trajectory~\cite{jastrzebski2020break}. In contrast, however, the Edge-of-Stability literature has shown that deep networks often continue to make progress near, or even beyond, the nominal sharpness threshold, suggesting that curvature alone is not a sufficient explanation~\cite{cohen2021edge, arora2022understanding, damian2023self}. Adaptive optimizer [adam, cite] adds additional complexity through memory and coordinate scaling: its adaptive preconditioner can produce loss spikes without an increasingly curved objective~\cite{bai2026spikes}.

These accounts motivate a further question about the learned function. Central Flows recognizes the connection between optimization dynamics and the roles of individual layers as an open problem~\cite{cohen2025central}. Neural loss landscapes are increasingly complex, multimodal, even fractal [add citations here, including the neurips visualization article and I think there are some in nature that are very strong]. The network therefore shouldn't be treated as a blackbox, and the optimization stability problem can actually be connected and mapped into the representation learning problem [this is us making an argument].

%%---- stop here
Do the responses inside the network change back and forth, and by how much? We follow those responses on fixed inputs and test whether limiting their local changes improves subsequent training. The argument proceeds from update interactions to functional measurements and then to matched interventions.

For plain stochastic gradient descent (SGD), gradient magnitudes and directions determine the applied increments. We distinguish persistent opposition between consecutive updates from cancellation over the full trajectory, and relate their interaction to directional curvature. This connection uses the classical secant construction underlying Barzilai--Borwein step selection~\cite{barzilai1988two,tan2016bbsgd}. Under Adam, raw gradients and applied increments require separate interpretations. We therefore distinguish exact accounting of a completed step from empirical warnings about the next step.

We study multilayer perceptrons (MLPs), residual networks, and Transformers on image-classification and language-modeling tasks. Full-objective training isolates update dynamics; mini-batch continuations examine sampling. Under Adam, fixed reference examples reveal representation changes, and interventions share the same saved training state. We make four contributions:
\begin{itemize}
\item \textbf{A framework for constructive and destructive dynamics.}
An exact SGD displacement decomposition separates step magnitudes from interactions across time. The experiments show that longer-lag reinforcement can outweigh adjacent opposition while loss falls. Sampling is examined separately at fixed parameters and along evolving trajectories (Sections~\ref{sec:displacement}--\ref{sec:gd-experiments}).

\item \textbf{Curvature measurement and empirical loss warnings.}
For smooth, fixed objectives, full gradients and the applied increment recover curvature averaged over a completed step without additional derivative evaluations. We connect this secant measurement to alignment, sampling, and the quadratic GD loss criterion, then separately evaluate fixed warning rules before the next Adam update (Sections~\ref{subsec:curvature-diagnostic}, \ref{subsec:loss}, and~\ref{subsec:characterization}).

\item \textbf{Measurement of functional oscillation.}
On fixed inputs, we measure the direction, magnitude, retention, and temporal structure of representation changes. The selected windows reveal repeated reversal along concentrated output patterns and different participation across network sites (Section~\ref{subsec:tracing}).

\item \textbf{A test of representation displacement as an intervention target.}
Across three matched model--dataset cases, every tested local-cap quantile reduces mean training loss and accumulated upward loss variation relative to unchanged Adam while reversals persist. Comparisons with conventional controls expose endpoint and computational tradeoffs (Section~\ref{sec:vc-interventions}).
\end{itemize}

Section~\ref{sec:background} positions these contributions in related work. Sections~\ref{sec:displacement}--\ref{sec:gd-experiments} develop and test the GD/SGD framework; Sections~\ref{sec:adam}--\ref{sec:vc-interventions} connect Adam diagnostics to functional measurements and interventions. Section~\ref{sec:conclusion} discusses implications and limitations.
